\begin{abstract}
Training 3D Gaussian Splatting (3DGS) on memory-limited GPUs is difficult because densification changes both persistent model state and transient optimizer memory, while aggressive pruning can remove newly useful primitives before they accumulate sufficient evidence. We present \method, a spatial resource-allocation framework that regulates clone and split mutations under a calibrated peak-VRAM budget. Its central mechanism is pressure-triggered preservation: low-opacity or low-utility Gaussians remain eligible for optimization while measured allocator headroom is sufficient and are evicted only when admission requires space. A dynamic inflow governor varies mutation admission between 5K and 15K candidates, and a gain-per-byte allocator selects feasible clone and split operations using calibrated affine costs. On an 8,150.56-MiB RTX 5060 Laptop GPU, we evaluate 24 physical 30K-step runs over Bonsai, Garden, and Flowers. In a seed-0 2x2 factorial study, preservation has a $+10.14$ dB PSNR main effect; the inflow-by-lifecycle interaction is $-3.41$ dB. Across three Bonsai seeds, preservation improves the memory-bounded baseline by $11.09\pm2.34$ dB (paired $t(2)=8.21$, $p=0.0145$), but the small sample precludes broad inferential claims. Single-seed gains are $+6.25$ dB on Garden and $+2.94$ dB on Flowers. All 24 runs complete without observed OOM; sampled physical peaks range from 6,906 to 7,656 MiB. These results establish empirical boundedness under the tested workload, not a hardware-level memory guarantee.
\end{abstract}
